% !TEX root = ../../../tree_notes.tex
% =====================================================================
%  12.2  The Diameter: Two Sweeps and a DP Join   sec:c09-diameter
%  Exports: def:c09-diameter  lem:c09-sweepend  thm:c09-twosweep
%           thm:c09-dpjoin  lst:c09-sweep  lst:c09-dp
%           fig:c09-diampath  ex:c09-diam-run
%           prob:c09-diam  lst:c09-1245-soln   (practice: CSES 2081)
%  Both algorithms, one answer, on the same tree T. The per-QUERY distance
%  (many pairs, online) stays Chapter 16's (c11) -- this section is
%  one-shot: compute D once, submit it.
% =====================================================================
\section{The Diameter --- Two Sweeps, One DP, Same Answer}\label{sec:c09-diameter}

\begin{previouslybox}
\begin{enumerate}[nosep, leftmargin=1.6em]
  \item \thmref{c09-onepass} folded $\height(u)=\max_c(\height(c)+1)$ bottom-up. What
        shape of path does that $\max$ single out at $u$? (the longest one that starts at
        $u$ and goes \emph{down})
  \item BFS from a source pops vertices in nondecreasing distance (used in
        \probref{c09-deep}). Why is ``pop order $=$ distance order'' enough to find a
        \emph{farthest} node without sorting anything? (the last popped has max distance)
  \item Every two nodes of a tree have exactly one path between them (\thmref{c01-unique-path},
        Ch.~1). If a path's \emph{topmost} node (nearest the root) is $t$, where do the
        path's two ends live? (in two different child-subtrees of $t$, or one is $t$)
  \item A single DFS from node $4$ of $T$ reports max distance~$6$. Is $6$ the diameter?
        (no --- $D=7$; the gap is exactly this section's first theorem)
\end{enumerate}
\end{previouslybox}

\begin{topicintro}
The \term{diameter} is the tree's longest simple path --- ``the number of edges in the
longest path'' is a literal problem statement (LeetCode~1245, \probref{c09-diam}), and a
literal ingredient in a dozen later problems. Doing it naively is
$\Theta(n^2)$: $\binom{n}{2}$ endpoint pairs, each with a distance. This section kills
the quadratic in two independent ways, and developing both is the point --- they fail
differently and generalise differently.

The first is the \term{two-sweep} trick: walk once from \emph{anywhere} to the farthest
node, then walk again from \emph{there}; the second distance is $D$
(\thmref{c09-twosweep}). The engine is one lemma worth memorising in its own right:
\emph{a farthest node from anywhere is always an endpoint of some diameter}
(\lemref{c09-sweepend}). The second way is a DP that never mentions farthest-ness: root
the tree, and observe every path has a topmost node where it bends across two distinct
downward arms --- so $D$ is the best ``join of the two longest arms'' over all nodes
(\thmref{c09-dpjoin}), computable in the same DFS as \coderef{c09-dfs3} (the arms
\emph{are} the $\height$ table). The two methods agree on every input precisely because
the lemma and the join theorem are both true --- and we prove both.

One boundary this section does not cross: distance \emph{queries} --- ``what is $d(u,v)$
for these $10^5$ pairs?'' --- is LCA arithmetic, Chapter~16 (\chref{c11}). Here we compute
the \emph{one} number $D$ (and optionally the path) once, which needs neither lifting nor
tour arrays. \textbf{Prerequisites.} \secref{c09-measures} (the $\height$ fold, the
BFS/DFS stacks), \thmref{c01-unique-path} (uniqueness of paths --- every proof below
silently leans on it).

\textbf{After this section you can}
\begin{itemize}[nosep]
  \item state and prove the sweep-endpoint lemma, and run the two sweeps on any tree; \hfill\lv{3}
  \item compute $D$ inside a \coderef{c09-dfs3}-style fold and explain what each method
        costs when edges get weights; \hfill\lv{3}--\lv{4}
  \item recover the actual diameter \emph{path}, not just its length; \hfill\lv{4}
\end{itemize}
\end{topicintro}

\subsection{Intuition}

Stand anywhere in the tree and walk to the node that is furthest away. Intuition says
that node should be an ``extremity'' --- and extremities are exactly where the longest
path has its ends. The two-sweep trick is this intuition cashed twice: the first walk
finds \emph{an} extremity $a$; the second walk, starting from an extremity, must reach
the \emph{opposite} end of a longest path, because nothing is further from an extremity
than the far extremity.

The DP reads the same object from the top down instead. Root the tree and look at any
path: climb it to its highest node $t$ --- the \emph{peak}. Below $t$ the path splits
into two arms going into two different child-subtrees (one arm may be empty, ending at
$t$ itself). So every path is ``arm $+$ peak $+$ arm'', and the longest path with peak
$t$ is just $t$'s two longest downward arms added together. Scan all peaks, take the
max --- that is a $\max\max$ over the same fold that produced $\height$
(\thmref{c01-height-recursion}), one DFS again.

Why bother having two methods at all? Because they degrade differently (§Variations):
with positive weights the DP carries over word for word while the sweeps need their
distance pass upgraded; the DP also extends to questions the sweeps cannot answer at all
(``the longest path satisfying \ldots'' once the local join gets a predicate). Knowing
both is what lets you \emph{choose} under contest pressure.

\subsection{Definitions and the Sweep Lemma}

\begin{definition}{Diameter and diameter path}{c09-diameter}
The \term{diameter} $D$ of a tree $T$ is the maximum number of edges over all simple
paths in $T$; a path achieving $D$ is a \term{diametral path}; a node that is an end of
some diametral path is \term{peripheral}. A single-vertex tree has $D=0$; ``diameter in
nodes'' (some LC problems) is $D+1$. The diametral \emph{length} is unique; the
\emph{path} need not be (a star has one per leaf pair).
\end{definition}

\begin{lemma}{Farthest is peripheral}{c09-sweepend}
Let $s$ be any node of a tree and let $x$ be a farthest node from $s$. Then $x$ is
peripheral: some diametral path has $x$ as an end.
\end{lemma}
\begin{proofbox}
Fix any diametral path $P = a \rightsquigarrow b$, $d(a,b)=D$. Let $m$ be $P$'s unique
node closest to $s$ (the path $s\rightsquigarrow P$ meets $P$ at $m$), and let $u$ be the
analogous attachment point of $x$. Because $P$ is simple and $x$ is off it only at $u$:
$d(s,x) = d(s,m) + d(m,u) + d(u,x)$, and $d(m,a)+d(m,b) = D = d(u,a)+d(u,b)$.

Two cases by where $u$ sits on $P$. \emph{If $u$ lies between $m$ and $b$}: comparing
$x$ against $a$, maximality of $x$ gives $d(m,x) \ge d(m,a)$, i.e.\
$d(m,u)+d(u,x) \ge d(m,u)+d(u,a)$, so $d(u,x) \ge d(u,a)$ and then
\[
  d(x,b) \;=\; d(u,x)+d(u,b) \;\ge\; d(u,a)+d(u,b) \;=\; D .
\]
\emph{If $u$ lies between $a$ and $m$}: comparing $x$ against $b$ instead gives
$d(u,x) \ge d(u,b)$ and then $d(x,a) \ge d(u,b)+d(u,a) = D$. ($u=m$ satisfies both.)

Either way some node is at distance $\ge D$ from $x$; no distance exceeds $D$ since $D$
is the maximum, so equality holds and $x$ ends a diametral path.
\end{proofbox}

\begin{theorem}{Two sweeps find the diameter}{c09-twosweep}
Pick any start node $s$. Let $a$ be a farthest node from $s$, and $b$ a farthest node
from $a$. Then $d(a,b) = D$.
\end{theorem}
\begin{proofbox}
By \lemref{c09-sweepend}, $a$ is peripheral: some $b'$ has $d(a,b') = D$. The second
sweep chooses $b$ to maximise $d(a,\cdot)$, so $d(a,b) \ge d(a,b') = D$, while $D$ is
the largest distance in the tree. Hence $d(a,b) = D$. Note what the proof used and what
it did not: it never matters \emph{which} node $s$ is, nor which farthest $a$ when ties
exist --- every choice lands peripheral.
\end{proofbox}

\begin{lstlisting}[caption={The two sweeps: one BFS (or DFS) distance pass, called twice; also returns the far node for path recovery}, label={lst:c09-sweep}]
// adj = undirected list, n nodes. sweep(src) = {farthest node, dist, par[]}
struct Sweep { int far; int dist; vector<int> par; };
Sweep sweep(int src) {
    vector<int> dist(n, -1), par(n, -1);
    queue<int> q;  q.push(src);  dist[src] = 0;
    int far = src;
    while (!q.empty()) {                       // BFS: pops in nondecreasing dist
        int u = q.front();  q.pop();
        if (dist[u] > dist[far]) far = u;      // ties: keeps first popped, any is fine
        for (int v : adj[u]) if (dist[v] == -1) {
            dist[v] = dist[u] + 1;  par[v] = u;
            q.push(v);
        }
    }
    return {far, dist[far], par};
}
int diameter(int any) {
    Sweep A = sweep(any);                    // sweep 1: reach an extremity A.far
    Sweep B = sweep(A.far);                  // sweep 2: B.dist = D; B.par walks the path
    return B.dist;                           // B.far -> ... -> A.far is a diametral path
}
// Thm: works for ANY start node, and ANY farthest when distances tie.
\end{lstlisting}

\begin{warningbox}[Read \coderef{c09-sweep}'s last comment before ``improving'' it]
The one-liner \code{D = sweep(sweep(any).far).dist} recomputes the second sweep
twice --- harmless but wasteful; the two-\code{auto} form is what you ship. What is
\emph{not} harmless: calling \code{sweep} \textbf{once} and submitting \code{A.dist}.
That is $x$'s distance from your chosen start, the eccentricity of $s$'s far corner, not
$D$ --- on $T$ from node~$4$ it reports $6$, not $7$. Trap~1 below.
\end{warningbox}

\subsection{The DP: Join the Two Longest Arms}

Root wherever convenient and let $\mathit{down}(u)$ be the longest \emph{downward} arm
from $u$ --- which is literally $\height(u)$ of \defref{c09-measures}. For a node $t$,
let $\mathit{arm}_1(t) \ge \mathit{arm}_2(t) \ge 0$ be the two largest values of
$\height(c)+1$ over the children $c$ of $t$ (missing second arm $=0$, i.e.\ a path
ending at $t$ is allowed).

\begin{theorem}{Diameter by joining arms at the peak}{c09-dpjoin}
$D = \max_t \big(\mathit{arm}_1(t) + \mathit{arm}_2(t)\big)$, computable in one DFS in
the style of \coderef{c09-dfs3}.
\end{theorem}
\begin{proofbox}
($\le$) Let $P$ be any simple path with peak $t$ (the unique node of $P$ nearest the
root --- uniqueness uses \thmref{c01-unique-path}). $P$'s two pieces below $t$ go into
two \emph{distinct} child-subtrees of $t$ (one piece may be empty, $t$ itself being a
path end), so each piece is at most an arm, and $|P| \le \mathit{arm}_1(t)+\mathit{arm}_2(t)$
by definition of $\height$. ($\ge$) Conversely take the child $c_1$ realising
$\mathit{arm}_1(t)$ and a deepest node $v_1$ in $\subtree(c_1)$: the path $t \rightsquigarrow
v_1$ has $\mathit{arm}_1(t)$ edges; similarly for $c_2 \ne c_1$. The concatenation
$v_1 \rightsquigarrow t \rightsquigarrow v_2$ is a simple path (distinct subtrees) of
exactly $\mathit{arm}_1(t)+\mathit{arm}_2(t)$ edges --- so the join value is achieved by
a real path, and the max over $t$ is both attainable and an upper bound.
\end{proofbox}

\begin{lstlisting}[caption={DP diameter inside the measures fold: keep the two longest arms at each node}, label={lst:c09-dp}]
vector<int> hgt2;  int D = 0;                 // hgt2[u] = down(u) = height (edges)
void dfs_diam(int u, int p) {
    hgt2[u] = 0;
    int a1 = 0, a2 = 0;                       // two longest arms via children
    for (int v : adj[u]) if (v != p) {
        dfs_diam(v, u);
        int arm = hgt2[v] + 1;
        if (arm >= a1) { a2 = a1; a1 = arm; }
        else if (arm > a2) a2 = arm;
    }
    D = max(D, a1 + a2);                      // best path with PEAK u
    hgt2[u] = a1;
}
// call dfs_diam(r, -1) for any r; D is root-choice-independent (why: the peak
// argument in the proof of Thm "Diameter by joining arms" works under every rooting).  Path recovery: also keep
// which child gave a1/a2, then walk down both.
\end{lstlisting}

\begin{complexitybox}
Both methods: $\Theta(n)$ time, $\bigO{n}$ memory (queues/arrays); \coderef{c09-dp} adds
$\bigO{h}$ stack and should be run through the \coderef{c09-iter} discipline on chain-
prone inputs, while \coderef{c09-sweep} is a queue-based BFS --- immune by construction.
Sweeps win when you also want \emph{all} distances from the ends (the $\mathit{dist}[]$
arrays are free, and \secref{c09-radius} spends them); the DP wins the moment anything
about the path stops being pure length (weights, constraints on arms, counting variants
like CSES~2081 in the practice set).
\end{complexitybox}

\subsection{Worked Example}

\begin{example}{Both methods on the running tree $T$}{c09-diam-run}
\emph{Two sweeps, start $4$.} BFS from $4$: farthest is $10$ at distance $6$ (the
$4,2,1,3,7,9,10$ chain) --- note $6 \ne D$, this single number is \emph{not} the
diameter. Second sweep from $10$: farthest is $8$ at distance $7$. Output $D=7$
(\thmref{c09-twosweep}).\\
\emph{DP join, root $1$.} Arms from the $\height$ table of \figref{c09-triples}: at $1$:
via~$2$: $\height(2){+}1=3$; via~$3$: $\height(3){+}1=4$; join $3{+}4=7$. At $2$:
$1{+}2=3$; at $3$: $1{+}3=4$; at $7$: $2{+}0=2$; leaves $0$. Max $=7$ --- same answer,
two unrelated proofs (\thmref{c09-dpjoin}). The path: the join at $1$ uses the arm
through $2$ down to $8$ and the arm through $3$ down to $10$:
\[
  8 - 5 - 2 - 1 - 3 - 7 - 9 - 10 \qquad (7 \text{ edges}). \qquad\checkmark
\]
\end{example}

\begin{figure}[H]
\centering
\begin{tikzpicture}[diagram, scale=0.85, transform shape,
    tnode/.style={circle, draw=primaryblue, line width=0.9pt, fill=primaryblue!8,
                  minimum size=6.2mm, font=\small\bfseries, text=primaryblue!60!black}]
  \node[tnode] (n1) at (0,0) {$1$};
  \node[tnode] (n2) at (-3,-1.6) {$2$};
  \node[tnode] (n3) at ( 3,-1.6) {$3$};
  \node[tnode] (n4) at (-5.2,-3.2) {$4$};
  \node[tnode] (n5) at (-1.6,-3.2) {$5$};
  \node[tnode] (n6) at ( 1.6,-3.2) {$6$};
  \node[tnode] (n7) at ( 5.2,-3.2) {$7$};
  \node[tnode] (n8) at (-1.6,-4.8) {$8$};
  \node[tnode] (n9) at ( 5.2,-4.8) {$9$};
  \node[tnode] (n10) at ( 5.2,-6.4) {$10$};
  \draw[tedge] (n1)--(n3) (n2)--(n4) (n3)--(n6);
  \draw[tedgehl] (n1)--(n2) (n2)--(n5) (n5)--(n8)
                 (n3)--(n7) (n7)--(n9) (n9)--(n10);
  \node[annot, font=\scriptsize, text=accentorange!80!black, anchor=south]
    at (n1.north) {peak $t=1$: arms $3$ (via $2$) $+\;4$ (via $3$) $=7$};
  \node[annot, font=\scriptsize, anchor=east] at (n8.west) {end $8$};
  \node[annot, font=\scriptsize, anchor=west] at (n10.east) {end $10$};
\end{tikzpicture}
\caption{The diametral path of $T$ in orange: $8\leftrightarrow10$, $D=7$ edges. The
off-path fringe (node~$4$ and node~$6$) hangs \emph{from the middle} of the diameter ---
which is the picture behind \lemref{c09-sweepend}: starting the first sweep at the fringe
node~$4$ still finds a true extremity ($10$), because a branch off the middle can never
outrun the nearer diameter end. Node~$1$ is where the DP does the join.}
\label{fig:c09-diampath}
\end{figure}

\begin{tracebox}[Sweep 2 in full: BFS from $10$ over $T$]
dist[10]=0; pop 10 -> 9 (1); pop 9 -> 7 (2); pop 7 -> 3 (3); pop 3 -> 1, 6 (4);
pop 6; pop 1 -> 2 (5); pop 2 -> 4, 5 (6); pop 4; pop 5 -> 8 (7).
max = 8 at dist 7 => D = 7.  \checkmark\ \ \code{par[]} chain back from 8:
8,5,2,1,3,7,9,10 -- the path in \figref{c09-diampath}. \checkmark
\end{tracebox}

\subsection{Variations and Twists}

\begin{itemize}[nosep]
  \item \textbf{Weighted edges (positive).} The DP transfers verbatim with
        $\mathit{arm} = \height(c) + w(u,c)$; the sweeps survive too --- the deep fact ---
        but the distance pass must compute true distances (DFS accumulation is fine on a
        tree; there is a unique path, no Dijkstra needed). Weighted tree DP lives in
        Chapters~13--14 (\chref{c10}); negative weights break \emph{both} ``farthest is
        peripheral'' steps. \hfill\lv{4}
  \item \textbf{The path itself.} Second sweep's $\mathit{par}[]$ array walks back from
        the far end to the start --- \coderef{c09-sweep}'s last line; every
        ``output one longest route'' problem is this two-line extra. \hfill\lv{3}
  \item \textbf{Counting paths by length.} Replace the $\max$ over arms by a
        \emph{count} of arm pairs with sum in a window --- the peak-join geometry is
        unchanged; CSES~2081 (practice) makes you earn it. \hfill\lv{5}
  \item \textbf{All-pairs longest (eccentricity table).} The sweeps give the two ends;
        getting $\mathit{ecc}(v)$ for \emph{every} $v$ is one more trick --- two distance
        arrays and a $\max$ --- \secref{c09-radius} picks this thread up. \hfill\lv{4}
\end{itemize}

\subsection{Common Pitfalls}

\begin{warningbox}[Trap 1 -- submitting the first sweep's distance]
One sweep returns $\max_v d(s,v)$: the eccentricity \emph{of your start}, which equals
$D$ only when $s$ happens to be peripheral. $T$ from node~$4$: returns $6$; from node~$1$:
returns $4$; the diameter is $7$. If your answer depends on which node you called the
sweep from, you have implemented a different function. (Symptom: WA only on chain-shaped
or lopsided tests, AC on stars --- from a star's centre every distance is $1$ and the
bug hides.)
\end{warningbox}

\begin{warningbox}[Trap 2 -- forgetting the second sweep needs fresh state]
\coderef{c09-sweep} allocates $\mathit{dist}/\mathit{par}$ inside the call; hand-rolled
globals reused across sweeps must be \emph{reset}. A stale $\mathit{dist}[]$ makes sweep
two see every node already visited: it returns the start node with distance~$0$ and your
``diameter'' collapses. This is the same class of bug as reusing one DFS array across
test cases --- and the same fix: reset per case, or timestamp the arrays.
\end{warningbox}

\begin{warningbox}[Trap 3 -- node-count vs edge-count, again]
LC~1245 counts \emph{edges} (\defref{c09-diameter}); LC~543-style binary problems count
\emph{nodes} on the path ($D+1$); ``height of a tree after removing all leaves'' families
mix a third convention. Decide from a two-node example before coding: two joined nodes
answer $1$ (edges) or $2$ (nodes) --- and check which one the sample prints.
\end{warningbox}

\subsection{Problem Lab}

The section's anchor is LeetCode~1245, \probref{c09-diam}: an edge list in, a number
out --- the purest possible call to \thmref{c09-twosweep}. It is also the smallest
input on which you should write both methods and \emph{diff} their answers.

% ---------------------------------------------------------------------------
\solvedproblem{c09-diam}{Tree Diameter}{LeetCode 1245}
  {https://leetcode.com/problems/tree-diameter/}{LC Medium · L3}
  {two sweeps, BFS distance pass, edge-list adjacency, path recovery}

\begin{problemstatement}[Tree Diameter (premium-era; statement and samples verified against the live problem)]
Given an undirected tree of $n$ nodes labelled $0 \dots n{-}1$, supplied as
\texttt{edges} with \texttt{edges.length} $= n-1$ and each entry
\texttt{[a,b]} a unit-weight edge, return the \emph{number of edges} in the longest
path. Constraints: $1 \le n \le 10^4$.
\textbf{Example~1}: \texttt{edges = [[0,1],[0,2]]} $\Rightarrow$ output \texttt{2}
(longest path $1-0-2$). \textbf{Example~2}: \texttt{edges =
[[0,1],[1,2],[2,3],[1,4],[4,5]]} $\Rightarrow$ output \texttt{4} (path $3-2-1-4-5$).
\end{problemstatement}

\begin{attemptgate}[Attempt it yourself first -- \lv{3}, give yourself 15 minutes]
$n \le 10^4$ forbids $\binom{n}{2}$ pairs with a traversal each ($\approx 10^8$
edge-visits at worst). You have two linear tools in this section; pick either.
\peek{How many BFS passes does \thmref{c09-twosweep} need, and from where?}
\nudge{Two: from any node (label~$0$ is always present and free to pick), then from the
farthest found. Each pass is a plain BFS over the adjacency list you build in $O(n)$.}
\fullhint{Build \code{adj} from \code{edges}; run \code{sweep(0)} $\to a$; run
\code{sweep(a)} $\to \{b, \mathit{dist}\}$; answer \code{dist}. Single-node trees
($n=1$, \code{edges} empty) must print $0$ --- make sure a sweep from the lone node
returns distance~$0$ and your code never reads \code{edges[0]}.}
\end{attemptgate}

\solstep{Understanding the Problem}
Input is the raw unrooted object --- an edge list with no root, no parent pointers,
sometimes a single node. Output is one integer: longest path \emph{in edges}. Nothing
asks for the path, but recovering it would cost one parent array, so the implementation
keeps it anyway (it is the same code as \figref{c09-diampath}'s caption walk-back).
The trap vocabulary: which node starts the sweep (any), what one sweep prints
(eccentricity of the start, not $D$), nodes-vs-edges ($1$-$0$-$2$: answer $2$).

\solstep{Key Observations}
\begin{observation}
An edge list with $\mathit{adj}[\cdot]$ is exactly \secref{c09-measures}'s data shape;
no rooting decision is even needed for the sweep method --- BFS does the orientation
implicitly via $\mathit{dist}[] == -1$ (visited $=$ ``claimed by a parent'').
\end{observation}
\begin{observation}
By \thmref{c09-twosweep} the answer is independent of the start node, so label~$0$ is
the canonical choice; by \lemref{c09-sweepend} the intermediate far node $a$ is
guaranteed peripheral, ties included.
\end{observation}
\begin{observation}
$n \le 10^4$ means recursion would actually survive (chains of $10^4$ frames are fine),
but BFS queues the immunity for free --- prefer \coderef{c09-sweep}'s shape and sleep.
\end{observation}

\solstep{Think It Through With Me}
\begin{thinkbox}
\deadend{Try every pair of nodes, BFS each, keep the max.}
{$\Theta(n^2)$--$\Theta(n^3)$-ish: $5\cdot 10^7$ pairs each sweeping $\bigO{n}$ ---
TLE, and it re-derives distances the two-sweep already collapses to two passes.}
\deadend{Root at 0 and answer $\height(\text{root})$ --- the deepest node's depth.}
{That is the longest path \emph{pointing down from 0}; a path whose peak is below 0 but
which does not pass through 0 (e.g.\ $3-2-1-4-5$ rooted at~$0$: its peak is node~$1$,
length~$4$, while $\height(0)=3$) is missed --- classic single-root fallacy;
\thmref{c09-dpjoin} fixes exactly this by joining \emph{two} arms and scanning all peaks.}
\aha{Two sweeps (\thmref{c09-twosweep}) --- or the same DFS as
\coderef{c09-dfs3} with the \coderef{c09-dp} join. Two linear passes, any start,
proof is \lemref{c09-sweepend}.}
\end{thinkbox}

\solstep{Solution}
Build adjacency from the edge list. Sweep (BFS) from $0$ to get a far node $a$; sweep
from $a$ to get a far node $b$ and distance $d$; return $d$. Correctness is
\thmref{c09-twosweep} verbatim; the $n=1$ degenerate case is handled because a BFS from
the lone node reports $\mathit{far}=0$, $d=0$.

\solstep{Proof of Correctness}
\begin{proofbox}[The returned \textit{dist} equals $D$]
Pass one runs \code{sweep(0)} and takes $a$ = a farthest node from~$0$; by
\lemref{c09-sweepend} $a$ is peripheral, so $\max_v d(a,v) = D$ (there is $b'$ with
$d(a,b')=D$ and no pair exceeds $D$). Pass two computes exactly that maximum: BFS pops
vertices in nondecreasing distance from $a$ (each enqueue adds one edge to a popped
distance, FIFO), so the final value of \code{far} realises $\max_v d(a,v)$, and the
returned \code{dist} is $D$. $n=1$: both sweeps run over one vertex, distances $\{0\}$,
return $0 = D$ for the single-vertex tree (\defref{c09-diameter}).
\end{proofbox}

\solstep{Implementation}
\begin{lstlisting}[caption={LeetCode 1245 -- two BFS sweeps over the edge-list adjacency}, label={lst:c09-1245-soln}]
class Solution {
public:
    int treeDiameter(vector<vector<int>>& edges) {
        int n = (int)edges.size() + 1;
        if (n == 1) return 0;                       // lone node: D = 0
        vector<vector<int>> adj(n);
        for (auto& e : edges) { adj[e[0]].push_back(e[1]); adj[e[1]].push_back(e[0]); }
        auto sweep = [&](int src, vector<int>* parOut) {
            vector<int> dist(n, -1), par(n, -1);
            queue<int> q;  q.push(src);  dist[src] = 0;
            int far = src;
            while (!q.empty()) {
                int u = q.front();  q.pop();
                if (dist[u] > dist[far]) far = u;
                for (int v : adj[u]) if (dist[v] == -1) {
                    dist[v] = dist[u] + 1;  par[v] = u;  q.push(v);
                }
            }
            if (parOut) *parOut = move(par);
            return pair<int,int>{far, dist[far]};   // {farthest, its distance}
        };
        int a = sweep(0, nullptr).first;            // sweep 1: reach an extremity
        vector<int> par;
        auto [b, d] = sweep(a, &par);               // sweep 2: D, + par[] for the path
        // (path: walk par backwards from b to a -- not needed for the answer)
        return d;
    }
};
\end{lstlisting}

\solstep{Trace}
\begin{tracebox}[Official Example 2: edges [0-1,1-2,2-3,1-4,4-5] $\Rightarrow$ 4]
adj: 0:\{1\}, 1:\{0,2,4\}, 2:\{1,3\}, 3:\{2\}, 4:\{1,5\}, 5:\{4\}.\\
Sweep from 0: pop 0(d0), 1(d1), 2(d2), 4(d2), 3(d3), 5(d3); far updates 0->1->2->3
(dist 3; node 5 also at 3 but ties keep the first) => a = 3.\\
Sweep from 3: pop 3(0), 2(1), 1(2), 4(3), 0(3), 5(4); far: 2,1,4, then 5 at dist 4
=> b = 5, d = 4. Return \textbf{4} \checkmark\ (path 3-2-1-4-5, official's own
explanation).\\
Example 1 [[0,1],[0,2]]: sweep(0): far=1 (d1, tie keeps first); sweep(1): pops 1(0),
0(1), 2(2) => d = \textbf{2} \checkmark. n=1: early return 0 \checkmark.
\end{tracebox}

\solstep{Complexity}
\begin{complexitybox}
Time: adjacency build $\Theta(n)$ plus two sweeps $\Theta(n)$ each $= \Theta(n)$
overall --- at $n = 10^4$ this is $\approx 4\cdot 10^4$ edge-visits, instantaneous.
Memory: $\Theta(n)$ (\code{adj}, two arrays per sweep, freed between calls). Verdict:
\ac. The DP method (\coderef{c09-dp}) is the same budget with recursion; either passes,
and you should be able to write both --- the mixed set asks for one that only the DP
shape survives.
\end{complexitybox}

\solstep{Pitfalls and Variants}
\begin{warningbox}
Slips here are the section's traps with faces: one sweep submitted (Trap~1 --- you would
have printed $3$ for Example~2), stale \code{dist} between sweeps (Trap~2), and
\code{edges.size()} read without the $+1$ for the node count. \emph{Variants worth
doing:} output an actual longest path (walk \code{par} --- keep it from sweep two only);
LC~543's binary-tree twin (the join DP through 120 lines of \chref{c02}-land, same
$\mathit{arm}_1+\mathit{arm}_2$ identity); and CSES~2081 below, where $\max$ becomes
$\#$ and the real fight starts.
\end{warningbox}

\begin{insightbox}[Key takeaways]
\begin{enumerate}[nosep, leftmargin=1.6em]
  \item Farthest-from-anywhere is a diameter end (\lemref{c09-sweepend}) --- the entire
        two-sweep trick is this lemma used twice (\thmref{c09-twosweep}).
  \item Every path has a unique peak; $D$ is the best two-arm join over all nodes
        (\thmref{c09-dpjoin}) --- the same fold that produced $\height$ in 12.1, with one
        more accumulator.
  \item One sweep distance is a \emph{starting-node eccentricity}, a different number;
        the second sweep is not redundant, it \emph{is} the diameter.
\end{enumerate}
\end{insightbox}

\begin{hardeningbox}
\begin{itemize}[nosep, leftmargin=1.4em]
  \item \textbf{Diff the methods.} Run \coderef{c09-sweep} and \coderef{c09-dp} on $T$
        by hand to $D=7$ (\exref{c09-diam-run}), then on the star $K_{1,4}$ ($D=2$, many
        diametral paths --- \defref{c09-diameter}'s uniqueness warning).
  \item \textbf{Prove the fringe claim.} In \figref{c09-diampath}, show node~$4$'s branch
        cannot produce a longer path than $8\!\leftrightarrow\!10$: a branch attaching at
        path position $p$ has depth $\le \min(d(p,a), d(p,b))$ or $P$ was not a diameter.
  \item \textbf{Recover paths cold.} Write the 5-line \code{par}-walk that prints
        $8,5,2,1,3,7,9,10$ from sweep two of \exref{c09-diam-run}.
\end{itemize}
\end{hardeningbox}

\subsection{Practice Problems}

\begin{stuckbox}
\begin{itemize}[nosep, leftmargin=1.4em]
  \item Did you call the sweep function \emph{twice}, and return the \emph{second}
        distance?
  \item Does your code return $0$ (not crash) for a one-node tree with empty
        \code{edges}?
  \item If you used the DP: are \code{a1}/\code{a2} reset \emph{inside} the per-node
        frame, and is $D$ updated with the \emph{sum}, while $\mathit{down}$ takes only
        the \emph{max}?
\end{itemize}
\end{stuckbox}

\begin{practiceset}
  \practice{Fixed-Length Paths II}{CSES 2081}
    {https://cses.fi/problemset/task/2081}{CSES · L5}
    {Count node pairs whose path length lies in $[a,b]$, $n \le 2\cdot 10^{5}$. This is
    \thmref{c09-dpjoin}'s geometry with $\max$ replaced by a count: sweep the peak $t$
    over all nodes; at $t$, a path is an arm from one child combined with an arm from
    another (or from $t$ itself), so maintain, per processed child, a table
    \code{cnt[depth]} of open arms, and for each new child's \code{depth} add the range
    of \code{cnt} entries that land in $[a,b]$. Merging child tables small-to-large
    keeps it near-linear; alternatively solve it by centroid decomposition once you
    reach Chapter~24 (\chref{c18}) --- this problem is that chapter's hello-world, so
    first pass: an honest $\bigO{n^2}$ ``sweep from every node'' to pin the geometry.}
\end{practiceset}

\begin{recallbox}[10-second refresher -- 12.2 The Diameter]
$D$ = longest path in edges (\defref{c09-diameter}). Method 1: sweep anywhere $\to a$,
sweep from $a$ $\to$ second distance $=D$ --- because farthest-from-anywhere is a
diameter end (\lemref{c09-sweepend}, \thmref{c09-twosweep}, \coderef{c09-sweep}).
Method 2: root, and join the two longest downward arms at every peak: $D=\max_t
(\mathit{arm}_1+\mathit{arm}_2)$ (\thmref{c09-dpjoin}, \coderef{c09-dp}) --- the same
fold as 12.1's $\height$. Path recovery $=$ second sweep's \code{par[]} walk-back. On
$T$: from~$4$ sweep gives~$6$ (a trap), sweeps give $10 \to 8$, $D=7$, path
$8\!-\!5\!-\!2\!-\!1\!-\!3\!-\!7\!-\!9\!-\!10$.
\selftest{
  \item State \lemref{c09-sweepend} and give its two-line consequence for algorithm design.
  \item Why does one sweep from node $4$ of $T$ print $6$ and not the diameter?
  \item Sketch the ($\le$) and ($\ge$) halves of \thmref{c09-dpjoin}'s proof.
  \item Which method keeps working (i) with positive weights, (ii) when asked to
        \emph{count} paths of length in $[a,b]$?
  \item Write the diametral path of $T$ from memory and its length.
}
\end{recallbox}
